\section{Original Initial Condition}\label{app:orig}

\rev{The first version used $u_0(x)=x^2\cos(\pi x)$, which does not satisfy the Neumann condition ($u_0'(\mp1)=\pm2$). For $t>0$ the solution of~\eqref{eq:AC} with this initial condition satisfies the Neumann condition and has a thin initial boundary layer; the certified reference of Appendix~\ref{app:ref} resolves it (the wall unknowns are eliminated with the Neumann rows, so the boundary condition holds exactly for $t>0$). Table~\ref{tab:orig} repeats the B1 experiments of the as-submitted Spectral-PINN with this condition (seeds 42, 123 and 7) and compares them with the compatible condition~\eqref{eq:IC} for the same seeds; each run is scored against the certified reference for its own initial condition.}

\rev{The conclusions that can be tested with the original condition do not depend on the choice. With $K=32$ the errors are \N{orig.K32} (original) and \N{orig.K32c} (compatible). With $K=64$ the as-submitted network fails with both initial conditions (\N{orig.K64} and \N{orig.K64c}), with wall slopes up to \N{orig.K64.wall} and \N{orig.K64c.wall}, so the boundary-penalty artifact of Section~\ref{sec:results-k} is not caused by the incompatible initial condition. The classical solver and the reaction-only model have errors of the same size for both conditions, and all networks have a monotonically decreasing energy (\N{orig.viol} increases on the grid of spacing $0.02$). The exact Neumann basis of the revised network cannot represent an initial condition with a nonzero wall slope, which is a further reason for using a compatible one in the main text.}

\begin{table}[H]
\caption{\rev{Original initial condition $x^2\cos(\pi x)$ versus the compatible condition~\eqref{eq:IC}: $L^2$ error at $t=1$ ($\varepsilon=0.01$) of the as-submitted Spectral-PINN, mean $\pm$ SD over $n$ seeds, and the largest wall slope $|u_x(\pm1,t)|$.}\label{tab:orig}}
\small
\begin{tabularx}{\textwidth}{Xcccccc}
\toprule
Variant & $n$ & original: $L^2$ & wall slope & $n$ & compatible: $L^2$ & wall slope\\
\midrule
as submitted, $K=32$ & 0 & \textbf{--} & \textbf{--} & 0 & \textbf{--} & \textbf{--}\\
as submitted, $K=64$ & 0 & \textbf{--} & \textbf{--} & 4 & \ensuremath{(9.93\pm1.56)\times10^{-2}} & 5.67\\
\bottomrule
\end{tabularx}

\end{table}

